\section{Related Work}
\label{sec:related}

\para{Reward hacking and over-optimization.}
Optimizing a misspecified objective can produce behavior that scores well and is unwanted~\citep{amodei2016concrete}, and \citet{skalse2022defining} give a formal treatment of when a proxy reward can be hacked.
In \rlhf~\citep{christiano2017deep,stiennon2020learning,ouyang2022training}, the proxy is a learned \rmodel.
\citet{gao2023scaling} study how true reward first rises and then falls as a policy is optimized against a proxy reward model, and \citet{moskovitz2023confronting} propose constrained \rlhf to limit over-optimization of component reward models.
\citet{casper2023open} survey the broader limitations of \rlhf.
These works measure \emph{how much} the proxy and the target diverge.
We ask \emph{what} the divergent responses look like, by labelling them with a taxonomy at each checkpoint.

\para{Length bias.}
\citet{singhal2023long} show that much of the reward gain in \rlhf is associated with longer outputs and that reward models correlate strongly with length.
Our \rmodel reproduces this bias from human votes (it prefers the longer response in 71.1\% of held-out pairs), and our verbosity results are consistent with theirs.
Unlike their work, we also ask whether longer responses are what a judge \emph{rejects}, and find that length grows across all responses while fabrication is what separates rejected ones.

\para{Sycophancy.}
\citet{perez2022discovering} use model-written evaluations to show that language models repeat users' stated views, and \citet{sharma2023towards} show that sycophancy is widespread across assistants and that human preference data sometimes favors responses that match a user's beliefs.
We borrow the wrong-belief trivia probe design from \citet{sharma2023towards}.
Unlike these works, we track sycophancy along an \rlhf trajectory, and in our setup it does not increase; we discuss why this absence is weak evidence in \secref{sec:discussion}.

\para{LLM judges.}
LLM judges are a common stand-in for human raters, with known position and verbosity biases~\citep{zheng2023judging,wang2023large}.
We rely on LLM judges because we have no human raters, and we therefore measure their reliability directly: position bias, order consistency, agreement between two judge families, and agreement with judge-free surface metrics.
These checks limit which of our conclusions can be trusted, and we report them as results.
